\documentclass[10pt,a4paper]{article}
\usepackage[margin=18mm,headheight=15pt]{geometry}
\usepackage{fontspec}
\setmainfont{NotoSerifCJK-Regular.ttc}[Path=/usr/share/fonts/opentype/noto/,FontIndex=2,BoldFont=NotoSerifCJK-Bold.ttc,ItalicFont=NotoSerifCJK-Regular.ttc,ItalicFeatures={FakeSlant=0.2}]
\setsansfont{NotoSansCJK-Regular.ttc}[Path=/usr/share/fonts/opentype/noto/,FontIndex=2,BoldFont=NotoSansCJK-Bold.ttc]
\setmonofont{DejaVu Sans Mono}
\XeTeXlinebreaklocale "zh"\XeTeXlinebreakskip=0pt plus 1pt
\usepackage{amsmath,amssymb,booktabs,array,tabularx,longtable,enumitem,graphicx,xcolor,tikz,fancyhdr}
\usetikzlibrary{arrows.meta,positioning,calc}
\usepackage{hyperref}
\hypersetup{hidelinks,pdftitle={高级图像处理·小测复习讲义（LaTeX黑白打印版）},pdfauthor={}}
\setlength{\parindent}{0pt}\setlength{\parskip}{4pt plus 1pt}
\linespread{1.12}\setlength{\emergencystretch}{3em}
\setlist{nosep,leftmargin=1.5em,topsep=3pt}
\renewcommand{\arraystretch}{1.25}
\pagestyle{fancy}\fancyhf{}\fancyhead[L]{\small 高级图像处理·小测复习讲义}\fancyhead[R]{\small 黑白打印版}\fancyfoot[C]{\thepage}
\begin{document}
\begin{center}{\LARGE\bfseries 高级图像处理·小测复习讲义}\\[4pt]{\large 黑白打印版 · LaTeX 重排}\end{center}
\vspace{3mm}
\section*{1. 模块一 · 绪论（Lecture01）}
\begin{itemize}
\item \textbf{三种定位}：图像处理是信号处理的一个专业领域；也是一种（工程）工具；知识链为 电路分析 $\to$ 信号与系统 $\to$ 随机信号处理 $\to$ 各类信号处理（图像处理）。
\item \textbf{目标的两种层次（必记）}：
\end{itemize}
\begin{center}\small\renewcommand{\arraystretch}{1.25}
\begin{tabularx}{\linewidth}{|>{\centering\arraybackslash}p{.19\linewidth}|>{\centering\arraybackslash}p{.13\linewidth}|X|p{.23\linewidth}|}\hline
\textbf{层次}&\textbf{英文}&\centering\textbf{面向}&\centering\arraybackslash\textbf{内容举例}\\\hline
低级图像处理\newline（Image improvement）&low level IP&面向\textbf{人的}观察与解释，改善图像的视觉效果&对比度拉伸、平滑、锐化、去噪等\\\hline
高级图像处理\newline（Image analysis）&high level IP&面向\textbf{机器的}自主感知，为测量特征与结构做准备&分割、特征提取、识别、压缩表示\\\hline
\end{tabularx}\end{center}
\textbf{本课程重点是低级 IP。}
\begin{itemize}
\item \textbf{易混概念}：图像处理 \& 计算机图形学（分析 vs 合成）；图像处理 \& 计算机视觉（低级 vs 高级阶段）；图像处理 \& 模式识别（如指纹识别）。图像识别是人工智能的一个应用领域，但两条技术线是\textbf{平行的}。
\item \textbf{图像获取设备}：摄像机（普通／鱼眼／全向）、红外相机、雷达、X 光机／CT、激光测距相机（Range camera）。
\item \textbf{处理链}：图像获取 $\to$ 图像增强 $\to$ 图像复原（去噪／几何校正／辐射校正）$\to$ 彩色图像处理 $\to$ 图像压缩（无损／有损）$\to$ 图像分割 $\to$ 表示与描述 $\to$ 识别。
\item \textbf{应用案例（能举 2～3 个即可）}：门禁／客流统计、A-B 门防尾随、遗留异物检测、车牌识别、虹桥机场飞机起降跟踪、无人机航拍交通参数与竣工图测绘。
\end{itemize}
\section*{2. 模块二 · 人眼视觉与感知（Lecture02-1）}
\begin{itemize}
\item \textbf{人眼结构}：角膜 + 巩膜（外膜）、脉络膜、睫状体、虹膜（隔膜）、晶状体、视网膜（含两类感光细胞）、盲点。
\item \textbf{两类感光细胞（必考对比）}：
\end{itemize}
\begin{center}\small\renewcommand{\arraystretch}{1.25}
\begin{tabularx}{\linewidth}{|>{\centering\arraybackslash}p{.16\linewidth}|p{.22\linewidth}|X|}\hline
\textbf{细胞}&\centering\textbf{视觉类型}&\centering\arraybackslash\textbf{分布与特性}\\\hline
\textbf{视锥细胞}\newline Cones&明视觉 / 亮视觉\newline (photopic)&集中于\textbf{中央凹}（fovea）；对\textbf{颜色}高度敏感；分辨细节能力强；在昏暗环境下不工作\\\hline
\textbf{视杆细胞}\newline Rods&暗视觉 (scotopic)&分布较广，提供总体视野（general view）；\textbf{无色觉}；暗光下灵敏、分辨力低\\\hline
\end{tabularx}\end{center}
\begin{itemize}
\item \textbf{成像}：晶状体是\textbf{柔性的}（区别于普通光学透镜，由睫状体纤维张力控制）；看远处变扁、看近处变厚；对应近视与远视。
\item \textbf{亮度适应（brightness adaptation）}：人眼总动态范围约 $10^{-6}\sim 10^4\ \mathrm{mL}$（毫朗伯），但\textbf{不能同时覆盖}；当前所处的灵敏度水平称为\textbf{亮度适应级}（光信号量化的依据）。
\item \textbf{韦伯比（Weber ratio）——高频考点}：令 $\Delta I_{50}$ 为对给定背景亮度 $I$ 时，有 $50\%$ 概率被回答“是”的照度增量，则 $\Delta I_{50}/I$ 称为韦伯比。\par
比值小 $\to$ 能区分很小的强度变化 $\to$ \textbf{亮度分辨能力好}；\textbf{比值大} $\to$ 只能区分较大的变化 $\to$ 分辨能力差。（背景越亮，能察觉的最小增量越大。）
\item \textbf{心理视觉效应（psycho-visual effects）——感知亮度不是强度的简单函数}：
\begin{itemize}
\item \textbf{马赫带（Mach band）}：在灰度渐变带的边界处，人眼会看到并不存在的“更亮/更暗”的条带，说明视觉系统对边缘有增强作用（侧抑制）。
\item \textbf{同时对比（simultaneous contrast）}：同一灰度的小块，放在暗背景上显得亮、放在亮背景上显得暗——感知亮度依赖于周围亮度。
\item \textbf{视错觉（optical illusion）}：如长度、方向、形状的误判。
\end{itemize}
\end{itemize}
\begin{quote}含义：\textbf{图像是为“人眼”增强的}，增强效果的好坏最终由人眼的主观感知决定——这也解释了为什么“客观均匀的直方图”未必等于“看起来最好”。\end{quote}
\section*{3. 模块三 · 采样、量化与像素关系（Lecture02-2）}
\begin{itemize}
\item \textbf{采样 vs 量化（必背）}：
\begin{itemize}
\item \textbf{采样（Sampling）}：在\textbf{空间}上离散化——把连续的 $(x,y)$ 变成网格点。图像的数学描述 $I=E(x,y,z,\lambda,t)$。
\item \textbf{量化（Quantization）}：在\textbf{幅度（灰度）}上离散化——把连续灰度变成有限级数。
\end{itemize}
两者合称\textbf{数字化}。
\item \textbf{采样定理（Nyquist）}：为不失真恢复，采样频率须 $f_s\geq 2f_{\max}$（$f_{\max}$ 为信号最高频率）。图像中含义：\textbf{采样间隔要足够小}，否则出现混叠（aliasing）——细密纹理变成虚假的低频条纹。
\item \textbf{两种分辨率（易混）}：
\begin{itemize}
\item \textbf{空间分辨率}：图像大小（行数 $\times$ 每行像素数）；同一范围内像素越多，分辨率越高。
\item \textbf{灰度分辨率}：每像素的比特数（bit）；位数越多，可区分的灰度级越多。
\end{itemize}
\item \textbf{分辨率的工程含义}：像素有物理尺寸 $\to$ 决定“图像中一个像素代表多大的实际尺寸”。\par
视场 $\mathrm{FOV}=\text{像素数}\times\text{像素尺寸}\times(\text{物距}/\text{焦距})$。示例：7952 px、像素 $4.52\ \mu\mathrm{m}$、物距 300 m、焦距 36.18 mm $\to$ 视场约 299 m，即\textbf{一个像素代表约 4.52 cm}。
\item \textbf{数字图像表示}：二值图像（0/1）／灰度图像（白 255、黑 0）／彩色图像（RGB，每个像素由三原色的灰度级描述）。
\item \textbf{像素关系（高频考点）}：
\begin{itemize}
\item \textbf{邻域}：4 邻域 $N_4(p)=$ 上下左右；对角邻域 $N_D(p)$；8 邻域 $N_8=N_4\cup N_D$。
\item \textbf{邻接}：4 邻接、8 邻接、\textbf{m 邻接（混合邻接）}——在 8 邻接基础上，\textbf{排除会造成“歧义多路连接”的对角邻接}，从而消除 8 邻接的路径二义性。
\item \textbf{路径}：4 路径 / 8 路径 / m 路径（连续像素的序列）。
\item \textbf{连通集}：只含一个连通分量；\textbf{区域}是连通集；\textbf{边界}：区域中至少有一个邻点不在该区域内的像素集合。
\end{itemize}
\item \textbf{算术运算}（逐像素）：加、减、乘、除。
\begin{itemize}
\item \textbf{图像平均（加 + 除）}：多帧相加取平均可去随机噪声，噪声方差降为 $\sigma^2/N$。
\item \textbf{图像相减}：提取差异／边缘／运动；$g=f-h$（如原图减模糊图或前一帧）。
\end{itemize}
\item \textbf{逻辑运算}：AND、OR、NOT（补）、XOR —— 常用于模板掩膜、二值图像处理与形态学。
\item \textbf{缩放与插值}：创建新像素位置 + 赋灰度；整数倍放大可用\textbf{像素复制}；插值方法：最近邻（nearest）、双线性（bilinear）、双三次（bicubic），平滑性与计算量递增。
\item \textbf{应用案例（10 m 气步枪报靶系统）}：环间距 10 mm、成绩需保留一位小数（0.1 环 = 1 mm）；需考虑①分辨率是否满足精度 ②如何实时判断子弹击中靶面 ③机位导致的透视变形如何校正 ④计算弹着点与靶心最近距离。
\end{itemize}
\section*{4. 模块四 · 空间域图像增强（重点 ★）}
\subsection*{4.1 基本框架}
\begin{itemize}
\item \textbf{增强 vs 复原}：\textbf{增强}是为特定应用让结果比原图“更适合”（problem-oriented，主观、无统一标准）；\textbf{复原}是利用对退化现象的\textbf{先验知识}恢复原图（客观）。
\item \textbf{三类处理途径}：①空间域\textbf{点处理}（只与单像素灰度有关，如对比度拉伸）；②空间域\textbf{掩模处理}（与邻域有关，如空间滤波）；③\textbf{频率域处理}（设计转移函数 $H(u,v)$，见第 5 节）。
\end{itemize}
\subsection*{4.2 灰度变换（点运算）}
\begin{center}\small\renewcommand{\arraystretch}{1.25}
\begin{tabularx}{\linewidth}{|>{\centering\arraybackslash}p{.16\linewidth}|p{.24\linewidth}|X|}\hline
\textbf{变换}&\centering\textbf{公式 / 参数}&\centering\arraybackslash\textbf{曲线特征与作用}\\\hline
\textbf{反色}\newline Image negative&$s=L-1-r$&负斜率直线；适合增强暗区（或亮区）中的白色／灰色细节，如数字减影、乳腺影像\\\hline
\textbf{对数}\newline Log transform&$s=c\cdot\log(1+r)$&上凸曲线；\textbf{扩展低（暗）灰度、压缩高灰度}，用于压缩动态范围（典型：显示傅里叶频谱）\\\hline
\textbf{幂律（伽马）}\newline Power-law / gamma&$s=c\cdot r^\gamma$&$\gamma<1\to$ 上凸，\textbf{提亮暗部、扩展低灰度}；$\gamma>1\to$ 下凹，压暗、扩展高灰度。\textbf{伽马校正}用于补偿显示/采集设备的非线性\\\hline
\textbf{分段线性}\newline Piecewise-linear&分段直线（需指定拐点）&形式可任意复杂（优点），但需要更多人工参数（缺点）；对比度拉伸是其特例\\\hline
\textbf{对比度拉伸}&$(r_1,s_1)$、$(r_2,s_2)$ 控制&把窄灰度范围拉宽到全动态范围，提高对比度\\\hline
\textbf{灰度级切片}\newline Gray-level slicing&阶跃 / 阶梯形&突出\textbf{特定灰度范围}（如缺陷、兴趣结构）；可只保留该段或保持其余不变\\\hline
\textbf{位平面切片}\newline Bit-plane slicing&取第 $k$ 个比特&把 8 bit 图分解为 8 个二值图；\textbf{高位}含主要视觉信息、\textbf{低位}含细节/噪声；用于分析各位的重要性、指导压缩与量化位数\\\hline
\end{tabularx}\end{center}
\subsection*{4.3 直方图处理}
\begin{itemize}
\item \textbf{定义}：$h(r_k)=n_k$（灰度为 $r_k$ 的像素数）；归一化（概率）$p(r_k)=n_k/n$，$k=0,\ldots,L-1$。
\item \textbf{读图规律}：整体偏左 $\to$ 图像偏暗；偏右 $\to$ 偏亮；分布范围宽 $\to$ 动态范围大、对比度高；集中在窄区间 $\to$ 对比度低；\textbf{双峰} $\to$ 前景/背景分离。
\item \textbf{直方图均衡（HE）}：以 \textbf{CDF 作为变换函数}，把任意分布映射为近似均匀分布。
\begin{itemize}
\item 连续：$s=T(r)=\int_0^r p_r(w)\,\mathrm{d}w$，且 $p_s(s)=1$（均匀），与 $p_r$ 无关。
\item 离散：$s_k=T(r_k)=(L-1)\sum_{j=0}^k p_r(r_j)$。
\item \textbf{步骤}：①统计直方图 $H[k]$ ②求累积 $H_c$ ③$T[k]=\operatorname{round}\bigl(H_c[k]\cdot(L-1)/MN\bigr)$ ④按 $T$ 重映射输出灰度。
\end{itemize}
\item \textbf{为什么 HE 后直方图不是完全均匀的？}（常考讨论题）①离散灰度级不可分割，一个输出级只能“吃下”整数个像素；②映射是\textbf{多对一}的（灰度级总数不增加），无法把像素均摊到每一级；③像素总数固定，只能整体搬移。
\item \textbf{规定化 / 匹配（specification）}：三步——①先对原图做均衡 ②指定目标 PDF 并求其变换函数 ③做\textbf{逆变换}：$z=G^{-1}(s)=G^{-1}(T(r))$。与均衡化的区别：均衡化目标是“一律均匀”，规定化是“匹配到指定形状”。
\item \textbf{局部增强}：整体均衡可能忽略小区域，对每个局部（小窗口）做处理可突出局部细节。
\end{itemize}
\subsection*{4.4 代数运算}
\begin{itemize}
\item \textbf{图像平均（去噪）}：$g'(x,y)=\frac{1}{N}\sum g_i(x,y)$；若噪声不相关且零均值，则 $E[g']=f$、$\sigma_{g'}^2=\sigma_n^2/N$ —— $N$ 越大越接近原图。
\item \textbf{图像相减}：$g=f-h$，用于提取差异/运动/边缘。
\end{itemize}
\subsection*{4.5 空间滤波}
\begin{itemize}
\item \textbf{掩模与相关/卷积}：$g(x,y)=\sum_{s=-a}^{a}\sum_{t=-b}^{b}w(s,t)\cdot f(x+s,y+t)$。步骤：选像素 $\to$ 取邻域 $\to$ 对邻域施加函数（返回标量）$\to$ 写回中心位置 $\to$ 遍历全图。
\item \textbf{边界延拓（padding）}：X（界外取常数 X）、Symmetric（镜像反射）、Replicate（复制最近边界值）、Circular（周期延拓）。
\item \textbf{平滑（低通）}：
\begin{itemize}
\item \textbf{均值滤波}：邻域平均；\textbf{系数和 = 1}（保持整体亮度）；可改尺寸；\textbf{缺点：模糊边缘与锐利特征}。
\item \textbf{中值滤波}：\textbf{非线性}，用邻域灰度中值替代中心；对\textbf{椒盐（脉冲）噪声}极有效，且比同尺寸线性平滑少模糊；但可能删除细小目标/纹理，对高斯噪声优势不明显。
\item \textbf{顺序统计滤波}：中值（median）、最大值（Max，找最亮点/膨胀）、最小值（Min，找最暗点/腐蚀）。
\end{itemize}
\item \textbf{锐化（高通）}：
\begin{itemize}
\item \textbf{一阶导数（梯度）}：产生\textbf{较粗的边缘}，对\textbf{灰度阶跃}响应强。
\item \textbf{二阶导数（拉普拉斯）}：对\textbf{细线、孤立点}响应更强，在阶跃处产生\textbf{双响应}（正负各一次）。
\item \textbf{拉普拉斯掩模}（4 邻域，系数和 = 0）：$\begin{bmatrix}0&-1&0\\-1&4&-1\\0&-1&0\end{bmatrix}$；锐化：$g=f-\nabla^2f$（或 $g=f+\nabla^2f$，取决于掩模符号约定）。
\item \textbf{非锐化掩模（unsharp masking）}：$\mathrm{highpass}=\text{原图}-\text{低通}$。
\item \textbf{高提升滤波（high-boost）}：$g=(A)\cdot\text{原图}-\text{低通}=(A-1)\cdot\text{原图}+\mathrm{highpass}$，$A>1$；视觉上“在锐化的同时保持原图亮度”。
\end{itemize}
\item \textbf{视觉表现}：高通滤波后在\textbf{暗背景上呈现增强的亮边缘}；平滑滤波后边缘变模糊。
\item 统计：平滑模板系数和 = 1（保持平均亮度），纯高通/锐化模板系数和 = 0（均匀区输出 0）。
\end{itemize}
\subsection*{4.6 进阶：HE 的失效与改进谱系（对应课件“进阶”章节）}
\begin{itemize}
\item \textbf{HE 的三个失效模式}：①平坦区域噪声被拉成“麻点”（bin 间距被拉开，噪声被放大）；②\textbf{亮度均值漂移}（AMBE 变大，整图“变脸”）；③多峰直方图失配。
\item \textbf{改进谱系（规律：每一代只修一个毛病）}：
\end{itemize}
\begin{center}\small\renewcommand{\arraystretch}{1.25}
\begin{tabularx}{\linewidth}{|>{\centering\arraybackslash}p{.25\linewidth}|X|}\hline
\textbf{方法（年份）}&\centering\arraybackslash\textbf{修的问题}\\\hline
AHE (1987)&分块各自均衡，解决全局映射不适应局部\\\hline
\textbf{CLAHE (1994)}&块内直方图\textbf{裁剪（clip limit）+ 双线性插值}，解决 AHE 放大噪声（工程标配）\\\hline
BBHE / DSIHE (1997/99)&\textbf{按均值/中位数}把直方图分两半各自均衡，缓解亮度均值漂移\\\hline
RMSHE / MMBEBHE\newline (2003/08)&递归分裂 / 最小化亮度误差 AMBE\\\hline
WTHE (2009)&按出现频率加权并截断，抑制过度增强\\\hline
Retinex / LIME (2016)&物理分解 $I=L\cdot R$，压光照 $L$、保反射 $R$\\\hline
Zero-DCE (2020)&网络直接学习灰度映射曲线 $T(r)$ —— 点操作的深度版\\\hline
\end{tabularx}\end{center}
\begin{itemize}
\item \textbf{直方图匹配 = 一维最优传输}：把源分布搬到目标分布 $z=G^{-1}(T(r))$，即两个 CDF 的复合（Monge 问题的一维解）。
\item \textbf{BBHE 的均值保持条件}：以均值 $m$ 分裂后，精确保持均值的充要条件是 $P(r\leq m)=1-m/255$。
\end{itemize}
\section*{5. 模块五 · 频域（变换域）图像增强（Enhancement3）}
\begin{quote}空间域是“拿掩模在图上滑”，频域是“按频率决定留谁、压谁”。两者由\textbf{卷积定理}严格等价，不是两套理论，只是“改哪个频率”比“改哪个邻域”更直观。\end{quote}
\subsection*{5.1 三个步骤与卷积定理}
\begin{itemize}
\item \textbf{流程（五步）}：$f\cdot(-1)^{x+y}\to$ \textbf{DFT} 得 $F(u,v)\to$ \textbf{乘滤波器} $G=F\cdot H\to$ \textbf{IDFT} $\to$ 取实部、再乘 $(-1)^{x+y}$、裁剪到 $[0,L-1]$。$H$ 称\textbf{转移函数}。
\item \textbf{乘 $(-1)^{x+y}$ 的作用}：由平移性质把\textbf{直流分量搬到频谱中心}（等价 fftshift），这样 $H$ 才能按“到中心的距离”$D(u,v)$ 设计——与课件 $D_0$ 的定义一致。
\item \textbf{卷积定理}：$f*h\Longleftrightarrow F(u,v)\cdot H(u,v)$。\textbf{补零}：DFT 隐含周期性，$h$ 须补零到 $M\times N$，否则是\textbf{循环卷积}——图像右边内容“绕回”左边，边界出现鬼影。
\item \textbf{何时用频域}：小掩模（$3\times3$）空域更省 $O(MN\cdot k^2)$、纯实数、无补零；\textbf{大掩模}或需按频率精确选择（低通／带阻／陷波／同态）时用频域，$O(MN\log MN)$ 与核尺寸无关（课件 “great saving for a large filter”）。
\end{itemize}
\subsection*{5.2 频谱图怎么读（课件只有图，结论补在这里）}
\begin{itemize}
\item \textbf{显示}：$|F(u,v)|$ 动态范围极大，须做\textbf{对数变换} $s=c\cdot\log(1+|F|)$，否则只能看到中心一个白点。
\item \textbf{位置 = 频率}：中心是\textbf{直流}（平均亮度），越往外频率越高。\textbf{亮度 = 该频率分量的幅度}，不是“此处像素亮”。
\end{itemize}

% 原稿第 4--7 页；正文忠实保留，包括原稿的技术表述。
\begin{itemize}
\item \textbf{对应关系：}平坦区／渐变 $\to$ 能量聚在中心（低频）；边缘、细线、纹理、噪声 $\to$ 外围亮线亮环（高频）；\textbf{等间距周期条纹 $\to$ 一对关于中心对称的亮点}；横线 $\to$ 过中心的竖直亮线（$\perp$ 条纹走向）。
\item \textbf{相位}携带位置与结构信息，只保留幅度无法重建原图；图像\textbf{平移}只改相位、幅度谱不变，\textbf{旋转}则频谱随之旋转。
\end{itemize}

\subsection*{5.3 低通与高通}
\begingroup\small
\renewcommand{\arraystretch}{1.3}
\begin{tabularx}{\linewidth}{@{}p{.10\linewidth}p{.25\linewidth}p{.10\linewidth}p{.10\linewidth}X@{}}
\toprule
\textbf{类型} & $\boldsymbol{H(u,v)}$ & $\boldsymbol{D=D_0}$ \textbf{处} & \textbf{振铃} & \textbf{特点} \\
\midrule
\textbf{ILPF}\par 理想 & $\begin{cases}1,&D\le D_0\\0,&D>D_0\end{cases}$ & 突跳 & 最重 & 反变换是 2D sinc（jinc），旁瓣无限、含负瓣；\textbf{物理不可实现} \\
\textbf{BLPF}\par 巴特沃斯 & $\dfrac{1}{1+(D/D_0)^{2n}}$ & $0.5$ & $n$ 小$\approx$无\par $n$ 大渐显 & 无明确截断；$n$ 为阶数，越大越接近理想 \\
\textbf{GLPF}\par 高斯 & $\exp\!\left(-\dfrac{D^2}{2D_0^2}\right)$ & $0.607$ & 无 & 反变换仍是高斯（无旁瓣），最平滑、边缘最不“硬” \\
\bottomrule
\end{tabularx}
\endgroup
\begin{itemize}
\item \textbf{振铃（ringing）：}$H$ 是理想矩形 $\to$ $h$ 是 sinc 型（无限旁瓣、含负瓣）$\to$ 阶跃边缘被抹成一圈圈明暗环。\textbf{截止越陡，振铃越重：}ILPF $>$ BLPF（$n$ 大）$>$ BLPF（$n$ 小）$>$ GLPF（无）。折中常用 \textbf{BLPF，$n=2\sim3$}。
\item \textbf{$D_0$ 怎么定：按功率比}——使半径 $D_0$ 的圆内包含总能量 \textbf{$90\%\sim95\%$}；$D_0$ 越小保留频率越少、越模糊。
\item \textbf{高通：}$H_{\mathrm{hp}}=1-H_{\mathrm{lp}}$（IHPF／BHPF／GHPF 与低通一一对应）。直流被去掉 $\to$ \textbf{暗背景上的亮边缘}，常“加回原图”（高提升）。频域低通 $\approx$ 空域均值／高斯，频域高通 $\approx$ 空域拉普拉斯／梯度；高通同样放大噪声。
\end{itemize}

\subsection*{5.4 带通／带阻与同态滤波}
\begin{itemize}
\item \textbf{带阻（bandreject）：}抑制以 $D_0$ 为中心、宽 $W$ 的环带（$|D-D_0|\le W/2$），其余通过——专治\textbf{周期性干扰}（电源纹波、网点、栅栏）。三种形式：理想、巴特沃斯
\[
H=\frac{1}{1+\left[DW/(D^2-D_0^2)\right]^{2n}}
\]
、高斯。\textbf{带通} $H_{\mathrm{bp}}=1-H_{\mathrm{br}}$，只保留该环带（提取特定尺度纹理）。
\item \textbf{陷波（notch）：}带阻的“点”版本，只挖一对特定频率。因实图像频谱\textbf{共轭对称}，噪声点必成对出现，\textbf{必须成对挖}，只挖一个会破坏对称性、反变换出虚部。课件 p13／p16 的“notch filter”指 $H(M/2,N/2)=0$，即只挖直流。
\item \textbf{同态滤波：}模型 $f=i\cdot r$——照明 $i$ 慢变$\to$\textbf{低频}，反射 $r$ 突变$\to$\textbf{高频}。\textbf{先取对数}把乘积变相加，两项才能在频域分别加增益：
\[
\ln\ \to\ \mathrm{DFT}\ \to\ \times H\ \to\ \mathrm{IDFT}\ \to\ \exp,
\qquad
H=(\gamma_H-\gamma_L)\bigl[1-\exp(-cD^2/D_0^2)\bigr]+\gamma_L.
\]
$\gamma_L<1$ \textbf{压低频}（压动态范围、去光照不均），$\gamma_H>1$ \textbf{提高频}（提局部对比度与细节）；$\gamma_L=\gamma_H$ 则退化为普通线性滤波。它既非单纯低通也非单纯高通。
\end{itemize}

\section*{6. 例题与标准答案（题库）}
带【练习1原题】标记的是课堂练习1 的原题，答案已写成可直接背诵的版本。

\noindent\textbf{题型索引：}A 掩模／滤波核特性判定 ｜ B 读图分析（直方图／频谱）＋方法选择 ｜ C 直方图／变换曲线作图 ｜ D 滤波原理问答 ｜ E 计算题 ｜ F 概念简答 ｜ G 应用题／开放设计 ｜ H 选择／判断小题

\subsection*{题型 A\quad 掩模／滤波核特性判定}
\subsubsection*{A1【练习1原题】指出下列掩模作用于图像，会产生什么特性。}
\[
\begin{array}{@{}ll@{\qquad\qquad}ll@{}}
\mathrm{a.}&\begin{bmatrix}2&-1&-1\\-1&2&-1\\-1&-1&2\end{bmatrix}&
\mathrm{b.}&\begin{bmatrix}0&-1&0\\-1&4&-1\\0&-1&0\end{bmatrix}\\[1.4em]
\mathrm{c.}&\dfrac12\begin{bmatrix}-2&1&-2\\1&6&1\\-2&1&-2\end{bmatrix}&
\mathrm{d.}&\dfrac19\begin{bmatrix}1&1&1\\1&1&1\\1&1&1\end{bmatrix}
\end{array}
\]
\textbf{答案：}
\begin{enumerate}[label=\alph*.,leftmargin=*]
\item 系数和 $=2-1-1-1+2-1-1-1+2=0$，是\textbf{二阶差分（方向性高通）}掩模。均匀区域响应接近 0；对\textbf{沿对角线方向}的灰度变化和边缘响应强，突出高频细节。
\item 系数和 $=4-1-1-1-1=0$，是\textbf{4 邻域拉普拉斯高通掩模}。突出\textbf{水平、垂直方向}的边缘、细线和孤立点，可用于锐化。
\item 系数和 $=\frac12\times(-2+1-2+1+6+1-2+1-2)=\frac12\times2=1$，中心系数较大，属于\textbf{保留平均亮度的锐化/高提升类掩模}（$\approx$ 原图 + 高频）。能增强中高频细节与边缘，同时保留直流分量（背景亮度不改变）。
\item 系数和 $=9\times(1/9)=1$，是 \textbf{$3\times3$ 均值（低通）掩模}。对邻域求平均，可平滑、降低随机噪声，但会\textbf{模糊边缘、细线和纹理}。
\end{enumerate}

\subsubsection*{A2\quad 判断下列三个掩模各属于低通还是高通、平滑还是锐化、是否有方向性。}
\[
\begin{array}{@{}ll@{\quad}ll@{\quad}ll@{}}
\text{①}&\begin{bmatrix}0&-1&0\\-1&5&-1\\0&-1&0\end{bmatrix}&
\text{②}&\begin{bmatrix}-1&-1&-1\\-1&8&-1\\-1&-1&-1\end{bmatrix}&
\text{③}&\begin{bmatrix}-1&-1&-1\\2&2&2\\-1&-1&-1\end{bmatrix}
\end{array}
\]
\textbf{答案：}

① 系数和 $=5-1-1-1-1=1\to$ \textbf{低通成分保留 + 拉普拉斯锐化}（即 $f+\nabla^2f$ 的等效核），是高提升/锐化掩模；\textbf{各向同性}（4 邻域拉普拉斯），无方向性。用于在保持亮度的同时增强边缘。

② 系数和 $=8-8=0\to$ \textbf{纯高通（八邻域拉普拉斯）}；各向同性，对\textbf{全方向}边缘与孤立点都敏感，输出在均匀区为 0。

③ 系数和 $=2+2+2-1\times6=0\to$ \textbf{有方向性的高通（一阶差分/类 Prewitt）}；上下两行系数为负、中间行为正，对\textbf{沿垂直方向的灰度变化（即水平走向的边缘）}响应最强，对垂直边缘不敏感。

\subsection*{题型 B\quad 读图分析（直方图／频谱）+ 增强／滤波方法选择}
\subsubsection*{B1【练习1原题】就题目中所给的 a、b、c、d 四个直方图，分析比较四幅图像的明暗状况和对比度高低。哪些直方图对应的图像需要处理以取得较好的可视效果？方法是什么？}
\textbf{通用规律：}直方图整体偏左 = 偏暗；偏右 = 偏亮；灰度分布范围越宽，动态范围与对比度越大；集中在窄区间 = 对比度低。

\begingroup\small\renewcommand{\arraystretch}{1.25}
\begin{tabularx}{\linewidth}{@{}cX@{}}
\toprule
\textbf{直方图} & \textbf{结论与处理方法}\\\midrule
a & 像素集中在\textbf{左侧暗区}、分布范围偏窄 $\to$ \textbf{整体偏暗、对比度低}。建议先做\textbf{对比度拉伸}；若还要整体提亮，用 $\gamma<1$ 的\textbf{伽马变换}或灰度平移。\\
b & 像素集中在\textbf{右侧亮区}、偏窄 $\to$ \textbf{整体偏亮、亮部层次不足}。可用 $\gamma>1$ 的\textbf{伽马变换}压暗，再做温和的对比度拉伸；已饱和的亮部细节无法恢复。\\
c & 像素集中在\textbf{中间窄区间} $\to$ 平均亮度适中但\textbf{对比度低}。建议\textbf{分段线性对比度拉伸}；需自动增强可用 \textbf{HE}，若有噪声或照明不均则用 \textbf{CLAHE}。\\
d & 灰度分布范围\textbf{接近全范围} $\to$ 动态范围和对比度都较好，\textbf{一般无需全局增强}。如仅有噪声，用中值滤波等保边方法。\\\bottomrule
\end{tabularx}\endgroup

\textbf{结论：}需要处理的是 \textbf{a、b、c}；d 一般不必再做全局增强。

\subsubsection*{B2\quad 一幅夜间道路监控图像，其直方图呈“两个高耸的峰，中间接近平坦”（双峰）。请说明它反映的图像特征，并给出增强方案。}
\textbf{答案：}双峰说明图像中大致存在\textbf{两类灰度占主导}的像素——一个峰对应暗背景（夜晚路面/天空），另一个峰对应较亮的车灯或反光目标；中间平坦表示中间灰度层次少，\textbf{目标与背景之间的过渡细节缺失}。

\textbf{增强思路：}①若目标就是“看清目标细节”，优先用 \textbf{CLAHE}（局部、限制对比度，不会把噪声放大），或对暗峰区域做 $\gamma<1$ + \textbf{分段线性拉伸}；②\textbf{慎用全局 HE}——两个峰会被强行拉开、平坦区被“拉出”大量噪声，亮峰区域可能过曝，反而降低可视效果；③若需保留场景真实感，先做噪声抑制（如中值/双边）再局部增强。

\subsubsection*{B3\quad 讨论：为什么直方图均衡化以后，变换图像的直方图往往并不是完全均匀的？}
\textbf{答案：}①连续性假设失效——“均匀分布”的结论是在灰度连续、像素无穷多的理想条件下得到的；

②\textbf{灰度级不可再分}：离散情况下输出灰度级数不超过输入级数，一个输出级只能接收整数个像素，无法均摊；

③\textbf{映射是多对一的}：多个输入灰度可能映射到同一个输出灰度，像素被“堆”在一起；

④像素总数固定，HE 只是把像素整体搬移，不能把某一级的像素“劈开”分配到多级上。

因此实际结果只是“比原来更接近均匀、对比度更高”，而不是严格均匀。

\subsubsection*{B4\quad 频谱图判读（Enhancement3）。(1) 中心亮点代表什么？(2) 远离中心的暗区代表什么？(3) 平坦的天空与密集的草地分别落在哪一部分？(4) 过中心出现一对对称亮斑，最可能是什么原因，用什么方法去除？}
\begin{enumerate}[label=(\arabic*),leftmargin=*]
\item \textbf{直流分量}（$u=v=0$），即图像\textbf{平均灰度}；多数图像能量集中在低频，故最亮。
\item \textbf{高频分量}，对应边缘、细线、纹理与噪声，能量小所以偏暗。
\item 天空（平坦、渐变）$\to$ \textbf{低频}，聚在中心附近；草地（细密纹理）$\to$ \textbf{高频}，外围亮带／亮环。
\item \textbf{周期性干扰}（电源纹波、网点、栅栏、扫描条纹）：单频正弦在频谱上是\textbf{一对关于中心对称的冲激}。用\textbf{带阻／陷波滤波}去除，离中心距离 = 条纹空间频率；条纹方向 $\perp$ 两点连线方向。
\end{enumerate}

\subsection*{题型 C\quad 直方图／变换曲线作图}
\subsubsection*{C1【练习1原题】画出题 3 图的直方图示意图，并给出增强操作的方法建议。（题图：左侧是较暗的人物与背景，右侧是一块较亮的投影屏幕）}
\textbf{手绘要点（不必精确复刻，画对坐标、峰位置与相对形状即可）：}
\begin{enumerate}[label=\arabic*.,leftmargin=*]
\item \textbf{画坐标轴：}横轴标“灰度级 $r$”，左端 0（黑）、右端 255（白）；纵轴标“像素数 $h(r)$”，只表示多少，不必标精确刻度。
\item \textbf{左侧暗部：}0--40 画一组较宽的高峰——靠近 0 处一个小尖峰（近黑背景），约 15--30 再画较宽隆起（人物与衣服）。
\item \textbf{中间低谷：}约 40--145 曲线贴近横轴，表示中间灰度像素很少。
\item \textbf{右侧亮部：}约 145--166 处画一个\textbf{窄而高}的峰（峰顶约 158），标注“投影屏幕”——屏幕上亮度接近，所以峰很集中。
\item \textbf{收尾标注：}亮峰右侧迅速降回横轴；左标“人物/背景暗部”、右标“投影屏幕亮部”。
\end{enumerate}
\textbf{增强方法建议：}
\begin{enumerate}[label=\arabic*.,leftmargin=*]
\item 优先用 \textbf{CLAHE} 做局部对比度增强，并设置适中的 clip limit，避免暗部噪声被过度放大。
\item 若主要目的是看清人物，对暗部用 $\gamma<1$ 的\textbf{伽马变换或分段线性拉伸}；不要直接强行提亮全图，以免屏幕区域进一步饱和。
\item 若屏幕内容同样重要，可\textbf{按亮度分区处理后再合成}。全局 HE 会同时拉伸两个峰，造成屏幕过亮、暗部噪声增强，\textbf{不宜作首选}。
\end{enumerate}

\subsubsection*{C2\quad 画出反色、对数变换 $s=c\cdot\log(1+r)$、幂律变换中 $\gamma<1$ 与 $\gamma>1$ 的灰度变换曲线示意（横轴 $r$，纵轴 $s$，均在 $[0,L-1]$）。}
\textbf{答案：反色：}过 $(0,L-1)$ 与 $(L-1,0)$ 的\textbf{负斜率直线}。

\textbf{对数：}过原点附近、\textbf{上凸}（增长率递减）的曲线，左端陡、右端平——\textbf{扩展暗部、压缩亮部}。

$\gamma<1$：过 $(0,0)$、$(L-1,L-1)$ 的\textbf{上凸}曲线（位于 $y=x$ 上方），提亮低灰度。

$\gamma>1$：过同样两点但\textbf{下凹}（位于 $y=x$ 下方），压暗低灰度、扩展高灰度。

（画图关键：所有幂律曲线都过 $(0,0)$ 和 $(L-1,L-1)$，$\gamma=1$ 时退化为 $y=x$ 恒等直线。）

\subsection*{题型 D\quad 滤波原理问答}
\subsubsection*{D1【练习1原题】对中值滤波回答下列问题：①如何操作 ②作用是什么 ③是否可以用模板操作 ④影响算法效率的因素 ⑤试估计其频域特性。}
① \textbf{如何操作？}对每个像素，取以它为中心的邻域窗口 $W$（常用 $3\times3$、$5\times5$ 等\textbf{奇数}尺寸），把窗口内灰度值\textbf{排序}，取\textbf{中间位置}的值替换中心像素，然后沿图像逐点滑动；边界处需规定复制/对称/循环等延拓方式。可写为
\[
g(x,y)=\operatorname{median}\{f(x+s,y+t)\mid(s,t)\in W\}.
\]
② \textbf{作用是什么？}它是\textbf{非线性平滑滤波}，对\textbf{椒盐（脉冲）噪声}特别有效——孤立极大/极小值通常不会成为排序后的中值；与同尺寸均值滤波相比，\textbf{边缘与细线模糊更小}，但可能删除小目标和细纹理，对高斯噪声的平滑优势不如均值类方法。

③ \textbf{能否用模板操作？}可以用模板（滑动窗口）定义\textbf{邻域}，但\textbf{不能}用固定系数模板像均值滤波那样做线性卷积。中值滤波需要\textbf{排序并取次序统计量}，输出随输入数据的排序关系变化，因此不存在对所有输入都适用的固定线性卷积掩模。

④ \textbf{影响效率的因素？}主要有：\textbf{窗口大小与形状}（样本数 $K$ 越大，比较/排序越多）；\textbf{图像尺寸} $M\times N$；\textbf{实现方法}（逐点完整排序约 $O(K\log K)$，选择算法、滑动窗口复用、直方图维护、并行化可加速）；\textbf{边界处理、数据类型、内存访问}与 SIMD/GPU 等硬件条件。

⑤ \textbf{频域特性？}中值滤波是\textbf{非线性系统}，不能写成固定的 $G(u,v)=F(u,v)H(u,v)$，因此\textbf{不存在}与输入无关的唯一频率响应。定性上表现为\textbf{自适应低通}：在平坦区抑制孤立的高频脉冲噪声，在边缘附近尽量保留阶跃结构，因而比均值低通更保边；窗口越大，等效低通作用越强、但边缘和小目标损失也越大。若要数值估计，可输入不同空间频率的正弦/条纹图，测输出幅度与输入幅度之比，得到给定窗口与信号条件下的\textbf{等效增益曲线}（工程上常用正弦扫描法）。

\subsubsection*{D2\quad 比较均值滤波与中值滤波在噪声抑制上的差异，分别适用哪种噪声？}
\textbf{答案：}均值滤波按邻域\textbf{平均}（线性低通），会模糊边缘，对\textbf{高斯噪声}有效（多个独立噪声相加后方差下降），但对\textbf{椒盐/脉冲噪声}效果差——均值会把噪声点的幅值“摊”到邻域里。中值滤波取邻域灰度的\textbf{中值}（非线性），能排除孤立的极值点，对\textbf{椒盐（脉冲）噪声}抑制效果显著，且能\textbf{较好地保留边缘和锐角}；但对高斯噪声的平滑能力弱于均值滤波。总之：\textbf{高斯噪声$\to$均值；椒盐噪声$\to$中值}。

\subsubsection*{D3\quad 为什么平滑（低通）模板的系数和等于 1，而纯锐化（高通）模板的系数和等于 0？}
\textbf{答案：}对一个\textbf{常数（均匀）区域}，平滑滤波的输出应当等于该常数，才能保持图像整体亮度不变，因此 $\sum w=1$。纯锐化/高通滤波的目的是\textbf{只保留灰度变化}、去掉直流（背景）成分，所以常数区域必须输出为 0，因此 $\sum w=0$。这也解释了“高通结果在暗背景上出现亮边”的视觉效果——只有梯度和细节被保留下来。

\subsubsection*{D4\quad 写出拉普拉斯锐化的公式与典型掩模，并说明一阶导数与二阶导数在锐化中的区别。}
\textbf{拉普拉斯（二阶）：}
\[
\nabla^2 f=f(x+1,y)+f(x-1,y)+f(x,y+1)+f(x,y-1)-4f(x,y),
\]
对应掩模
\[
\begin{bmatrix}0&-1&0\\-1&4&-1\\0&-1&0\end{bmatrix}
\]
（还有含对角项的 8 邻域版本，中心系数为 8）。锐化输出 $g=f-\nabla^2f$ 或 $g=f+\nabla^2f$（由掩模符号约定决定）。

\textbf{一阶 vs 二阶：}①一阶导数（梯度）一般产生\textbf{较粗}的边缘，对\textbf{灰度阶跃}响应更强；②二阶导数对\textbf{细线、孤立点}等细节响应更强，且在阶跃处产生\textbf{双响应}（一正一负）；③因此增强细线细节偏好二阶（拉普拉斯），而检测边缘/轮廓常用一阶（梯度/Roberts/Sobel/Prewitt）。

\subsubsection*{D5\quad 说明非锐化掩模与高提升滤波的关系。}
\textbf{答案：非锐化掩模：}用原图减去其\textbf{模糊（低通）}版本得到高通（锐化）分量，即 $\mathrm{highpass}=f-f_{\mathrm{blur}}$。它丢失了原图的整体亮度。\textbf{高提升滤波}是在此基础上把原图按倍数 $A$ 加回：
\[
g=(A)\cdot f-f_{\mathrm{blur}}=(A-1)\cdot f+\mathrm{highpass},\qquad A>1.
\]
当 $A=1$ 时退化为非锐化掩模；$A$ 越大，原图（低频）占比越高，效果是\textbf{“在保持/提高整体亮度的同时锐化细节”}。

\subsubsection*{D6\quad 写出频域滤波的完整流程；并说明：(1) 为什么变换前要乘 $(-1)^{x+y}$？(2) 为什么补零到 $M\times N$？(3) 什么时候该用频域而不是 $3\times3$ 掩模？}
\textbf{流程：}
\[
f\cdot(-1)^{x+y}\ \to\ \mathrm{DFT}\ \to\ G=F\cdot H\ \to\ \mathrm{IDFT}
\]
$\to$ 取实部、再乘 $(-1)^{x+y}$、裁剪到 $[0,L-1]$。
\begin{enumerate}[label=(\arabic*),leftmargin=*]
\item 由平移性质，乘 $(-1)^{x+y}$ 把\textbf{直流搬到频谱中心}（等价 fftshift），$H$ 才能写成以中心为原点的 $D(u,v)$，否则要写成四个象限的副本。
\item DFT 隐含周期性，不补零就是\textbf{循环卷积}——图像右边内容绕回左边，边界出鬼影；补零到同尺寸后才是线性卷积。（$H$ 若直接在频域构造，如 ILPF，则无此问题。）
\item 小掩模（$3\times3$）用空域更省：$O(MN\cdot k^2)$、纯实数、无补零；\textbf{大掩模或需按频率精确选择}（带阻、陷波、同态）时用频域，$O(MN\log MN)$ 与核尺寸无关。
\end{enumerate}

\subsubsection*{D7\quad 同态滤波。(1) 写出模型，说明照明与反射各对应低频还是高频；(2) 为什么必须先取对数？(3) 写出流程与 $H(u,v)$，并说明 $\gamma_L$、$\gamma_H$ 的作用。}
\begin{enumerate}[label=(\arabic*),leftmargin=*]
\item $f=i\cdot r$，$0<i<\infty$，$0<r<1$；照明 $i$ 空间慢变 $\to$ \textbf{低频}，反射 $r$ 在物体交界处突变 $\to$ \textbf{高频}。
\item $f$ 是 $i$ 与 $r$ 的乘积，频域里乘积变卷积，无法“只动低频、不动高频”；取对数后 $z=\ln i+\ln r$ 变相加，两项可分离，同一个 $H$ 才能给不同增益。
\item $\ln\to\mathrm{DFT}\to\times H\to\mathrm{IDFT}\to\exp$；
\[
H=(\gamma_H-\gamma_L)\bigl[1-\exp(-cD^2/D_0^2)\bigr]+\gamma_L.
\]
$\gamma_L<1$ 压低频 $\to$ 压动态范围、去光照不均；$\gamma_H>1$ 提高频 $\to$ 提局部对比度。$\gamma_L=\gamma_H$ 时退化为普通线性滤波。
\end{enumerate}
\textbf{易错：}同态滤波既非单纯低通也非单纯高通，而是“低频压、高频提”——在 log 域线性，对原图是非线性。

% Faithful transcription of source pages 8--10. Source issues are recorded separately.
\clearpage
\section*{题型 E\quad 计算题}
\subsection*{E1\quad 对一幅 $64\times64$、8 个灰度级的图像做直方图均衡。各灰度级像素数如下，求均衡后的映射关系与新直方图。}
\begin{center}
\small\setlength{\tabcolsep}{4pt}
\begin{tabularx}{\linewidth}{|l|*{8}{>{\centering\arraybackslash}X|}}\hline
\textbf{灰度级 $r_k$}&0&1&2&3&4&5&6&7\\\hline
像素数 $n_k$&790&1023&850&656&329&245&122&81\\\hline
\end{tabularx}
\end{center}
\textbf{第 1 步\quad 求概率。}总像素 $MN=64\times64=4096$；$p(r_k)=n_k/4096$：0.19, 0.25, 0.21, 0.16, 0.08, 0.06, 0.03, 0.02（和为 1.00）。

\textbf{第 2 步\quad 求累积 CDF。}0.19, 0.44, 0.65, 0.81, 0.89, 0.95, 0.98, 1.00。

\textbf{第 3 步\quad 求输出灰度。}$s_k=\operatorname{round}(\mathrm{CDF}_k\times(L-1))=\operatorname{round}(\mathrm{CDF}_k\times7)$：

$0.19\times7=1.33\to1$；$0.44\times7=3.08\to3$；$0.65\times7=4.55\to5$；$0.81\times7=5.67\to6$；$0.89\times7=6.23\to6$；$0.95\times7=6.65\to7$；$0.98\times7=6.86\to7$；$1.00\times7=7\to7$。

\textbf{第 4 步\quad 得映射：}$0\to1$，$1\to3$，$2\to5$，$3\to6$，$4\to6$，$5\to7$，$6\to7$，$7\to7$。

\textbf{第 5 步\quad 合并像素、求新直方图。}
\begin{center}
\small\setlength{\tabcolsep}{4pt}
\begin{tabularx}{\linewidth}{|l|*{8}{>{\centering\arraybackslash}X|}}\hline
\textbf{新灰度 $s_k$}&0&1&2&3&4&5&6&7\\\hline
构成&—&$r=0$&—&$r=1$&—&$r=2$&$r=3,4$&$r=5,6,7$\\\hline
新像素数 $n_k$&0&790&0&1023&0&850&985&448\\\hline
新概率 $p(s_k)$&0&0.19&0&0.25&0&0.21&0.24&0.11\\\hline
\end{tabularx}
\end{center}
\textbf{说明：}新直方图比原来更均匀（原为 0.19/\allowbreak 0.25/\allowbreak 0.21/\allowbreak 0.16/\allowbreak 0.08/\allowbreak 0.06/\allowbreak 0.03/\allowbreak 0.02，现为 0.19/\allowbreak 0.25/\allowbreak 0.21/\allowbreak 0.24/\allowbreak 0.11），但不完全均匀——因为灰度级数没变、映射是多对一的。

\subsection*{E2\quad 分辨率的工程换算。某相机像素数 $7952\times5304$，像素尺寸 $4.52\,\mu\mathrm{m}$，物距 $300\,\mathrm{m}$，镜头焦距 $36.18\,\mathrm{mm}$。求：(1) 传感器尺寸；(2) 视场（横向 $\times$ 纵向）；(3) 图像中一个像素代表多大的实际尺寸。}
\textbf{(1) 传感器尺寸：}$7952\times4.52\,\mu\mathrm{m}\approx35.94\,\mathrm{mm}$（横）；$5304\times4.52\,\mu\mathrm{m}\approx23.97\,\mathrm{mm}$（纵）。

\textbf{(2) 视场：}由几何关系 $\mathrm{FOV}=\text{传感器尺寸}\times\text{物距}/\text{焦距}$：横 $=35.94\,\mathrm{mm}\times300\,\mathrm{m}/36.18\,\mathrm{mm}\approx299\,\mathrm{m}$；纵 $=23.97\times300/36.18\approx199\,\mathrm{m}$。

\textbf{(3) 单像素对应的物方尺寸：}$299\,\mathrm{m}/7952\approx0.0376\,\mathrm{m}\approx3.8\,\mathrm{cm}$（按纵向算 $199/5304\approx3.75\,\mathrm{cm}$）。题目原场景给出约 \textbf{4.52 cm/像素}，量级一致。

\textbf{结论：}该相机在 300 m 处分辨能力约 4 cm/像素，\textbf{无法分辨毫米级裂纹}——这正是隧道检测案例要求“更高分辨率/更近物距/多相机”的原因。

\subsection*{E3\quad 图像平均去噪。对同一静止场景连续拍摄 $N$ 帧并逐像素取平均，若各帧噪声互不相关、零均值、方差为 $\sigma^2$，求平均后图像的噪声方差与标准差；若要求噪声标准差降为原来的 $1/5$，需要平均多少帧？}
\textbf{答案：}$g'(x,y)=\dfrac{1}{N}\sum_{i=1}^{N}g_i(x,y)$。因噪声不相关且零均值：$E[g']=f$（无偏），$\sigma^2_{g'}=\sigma^2/N$，标准差 $\sigma_{g'}=\sigma/\sqrt{N}$。

要求 $\sigma/\sqrt{N}=\sigma/5\to\sqrt{N}=5\to N=25$ 帧。即噪声标准差按 $1/\sqrt{N}$ 下降，帧数要增加到平方倍。

\subsection*{E4\quad 韦伯比。(1) 写出韦伯比的定义并说明其含义；(2) 某背景亮度 $I=100$（单位任意），在 $50\%$ 概率下可察觉的最小照度增量 $\Delta I_{50}=2$，求韦伯比；(3) 换一个更亮的背景 $I=400$、$\Delta I_{50}=16$，比较两次的亮度分辨能力。}
\textbf{(1) 定义：}$\Delta I_{50}/I$，其中 $I$ 为背景亮度，$\Delta I_{50}$ 为有 $50\%$ 概率被回答“是”的照度增量。比值\textbf{小}表示能区分很小的强度变化 $\to$ 分辨能力\textbf{好}；比值\textbf{大}表示只能区分较大的变化 $\to$ 分辨能力\textbf{差}。

\textbf{(2)} 韦伯比 $=2/100=0.02$。

\textbf{(3)} 第二次 $=16/400=0.04$。两者韦伯比不同：第二次的比值更大，说明在更亮的背景下，人眼\textbf{分辨能力下降}——背景越亮，需更大的绝对增量才能被察觉（这正是“暗适应”比“亮适应”更灵敏的体现，也是量化级差设计要考虑韦伯定律的原因）。

\subsection*{E5\quad 采样定理。某图像最细结构的空间周期为 $0.2\,\mathrm{mm}$，要在数字化时不发生混叠，采样间隔最大为多少？采样频率至少为多少？}
\textbf{答案：}最细结构周期 $T=0.2\,\mathrm{mm}\to$ 最高空间频率 $f_{\max}=1/T=5\,\mathrm{cycles/mm}$。由采样定理 $f_s\geq2f_{\max}=10$ 次/mm，对应最大采样间隔 $\Delta x\leq1/f_s=0.1\,\mathrm{mm}$。

\textbf{含义：}像素间距（采样间隔）必须 $\leq0.1\,\mathrm{mm}$，即像素尺寸至少要比最细结构小一半，否则细密纹理会出现虚假的低频条带（混叠）。工程上一般再留 2～3 倍裕量，取 0.03–0.05 mm/像素。

\subsection*{E6\quad 频率域计算 + 振铃。$M=N=256$，频谱已中心化，$D_0=32$，某点偏移 $(\Delta u,\Delta v)=(24,32)$。(1) 求 $D$；(2) 求 ILPF、BLPF($n=2$)、GLPF 在该点的 $H$；(3) 用 ILPF 平滑含边界的图像，为什么边界两侧出现一圈圈明暗环？哪种滤波器不振铃？}
\textbf{(1)} $D=\sqrt{24^2+32^2}=\sqrt{1600}=40$。

\textbf{(2)} 因 $40>32$：ILPF $H=0$；BLPF $H=1/[1+(40/32)^4]=1/(1+2.441)=0.291$；GLPF $H=\exp(-40^2/(2\times32^2))=\exp(-0.781)=0.458$。

补充：$D=D_0$ 时 ILPF 处于跳变边界、BLPF 恒为 \textbf{0.5}（$-3\,\mathrm{dB}$，与 $n$ 无关）、GLPF 为 $e^{-0.5}\approx0.607$。

\textbf{(3) 即振铃：}ILPF 的 $H$ 是理想矩形，其反变换 $h$ 是 $\operatorname{sinc}$ 型函数，带\textbf{无限旁瓣（含负瓣）}；滤波 $=$ 与 $h$ 卷积，阶跃边缘被旁瓣抹成一圈圈明暗环。\textbf{截止越陡，振铃越重}（ILPF $>$ BLPF 高阶 $>$ BLPF 低阶 $>$ GLPF）。\textbf{GLPF 不振铃}，因为高斯的傅里叶变换仍是高斯、没有旁瓣；代价是截止不陡、边缘略糊。工程折中取 BLPF $n=2\sim3$。

\section*{题型 F\quad 概念简答}
\subsection*{F1\quad 图像增强与图像复原的区别是什么？}
\textbf{答案：图像增强}是为了让结果比原图\textbf{更适合某个特定应用}而进行处理，是\textbf{面向问题（problem-oriented）}的，如对比度拉伸、平滑、锐化；它\textbf{没有统一的客观标准}，效果好坏常由人眼主观判断。\textbf{图像复原}是利用对\textbf{退化现象的先验知识}，把被退化的图像恢复到接近原图，是\textbf{客观}的，追求“逼近真实原图”。一句话：增强是“让它更好看/更有用”，复原是“把它还原回去”。

\subsection*{F2\quad 低级图像处理与高级图像处理的区别？本课程的重点是哪一类？}
\textbf{答案：低级图像处理（Image improvement）}面向\textbf{人}的观察与解释，目标是改善图像的视觉效果，如增强、去噪、锐化；\textbf{高级图像处理（Image analysis）}面向\textbf{机器}的自主感知，为测量图像中的特征与结构做准备，如分割、特征提取、识别、压缩表示。\textbf{本课程重点是低级图像处理。}

\subsection*{F3\quad 视锥细胞与视杆细胞有什么差别？这对图像设计有什么意义？}
\textbf{答案：视锥细胞}集中在\textbf{中央凹}，负责\textbf{明视觉（亮视觉）}，对\textbf{颜色}高度敏感、分辨细节能力强，但暗光下不工作；\textbf{视杆细胞}分布更广，负责\textbf{暗视觉}，提供总体视野、\textbf{无色觉}，暗光下灵敏但分辨力低。

\textbf{意义：}人眼对亮度的感觉主要由视杆细胞完成，而其响应相当于一定范围内多个细胞的\textbf{平均}——这解释了为什么中值滤波等“不改变视觉感受”的去噪方法在主观上看不出损失；也说明图像增强的效果最终要以人眼主观感受为准（马赫带、同时对比等即为证据）。

\subsection*{F4\quad 什么是采样？什么是量化？图像数字化包含哪两方面？}
\textbf{答案：采样}是在\textbf{空间位置}上对连续图像离散化（把连续的坐标变成网格点）；\textbf{量化}是在\textbf{幅度（灰度）}上离散化（把连续灰度变成有限个灰度级）。图像的数学描述为 $I=E(x,y,z,\lambda,t)$；\textbf{数字化 $=$ 空间位置的数字化（采样）$+$ 灰度的数字化（量化）。}

\subsection*{F5\quad 空间分辨率与灰度分辨率的区别？分别由什么决定？}
\textbf{答案：空间分辨率}由图像中\textbf{像素的密集程度}决定（同一范围内像素越多分辨率越高，即行数$\times$列数）；\textbf{灰度分辨率}由每个像素的\textbf{比特数}决定（位数越多，可区分的灰度级越多，如 8 bit $\to$ 256 级）。降低空间分辨率会使图像出现块状/马赛克，降低灰度分辨率会出现“假轮廓”（伪轮廓）。

\subsection*{F6\quad 写出像素 $p$ 的 4 邻域、对角邻域、8 邻域；什么是 $m$ 邻接？为什么要引入 $m$ 邻接？}
\textbf{答案：}对像素 $p(i,j)$：\textbf{4 邻域} $N_4=\{(i-1,j),(i+1,j),(i,j-1),(i,j+1)\}$；\textbf{对角邻域} $N_D=\{(i-1,j-1),(i-1,j+1),(i+1,j-1),(i+1,j+1)\}$；\textbf{8 邻域} $N_8=N_4\cup N_D$。

\textbf{$m$ 邻接（混合邻接）：}在 8 邻接的基础上，\textbf{排除会造成“多路歧义连接”的对角邻接}——即若两像素既 4 邻接（或可通过 4 邻接到达），则不再承认其对角邻接。其意义是\textbf{消除 8 邻接在路径/连通判断上的二义性}，使连通路径唯一、更符合直观。

\subsection*{F7\quad 图像的算术运算与逻辑运算各有什么用途？各举一例。}
\textbf{答案：算术运算}（加/减/乘/除，逐像素）：\textbf{加法+平均}用于多帧平均去随机噪声（$\sigma^2\to\sigma^2/N$）；\textbf{减法}用于提取差异、运动或边缘（$g=f-h$）；乘法/除法常用于阴影校正、掩膜。\textbf{逻辑运算}（AND/OR/NOT/XOR）：用于二值图像处理、模板掩膜与形态学——如用 AND 提取感兴趣区域、XOR 找出两幅二值图的差异。

\section*{题型 G\quad 应用题／开放设计}
\subsection*{G1\quad （气步枪报靶）$10\,\mathrm{m}$ 气步枪靶环间距 $10\,\mathrm{mm}$，成绩需保留一位小数（0.1 环 $=1\,\mathrm{mm}$）。请说明分辨率如何满足精度要求，并简述其余三个工程问题。}
① \textbf{分辨率要求：}成绩精度 0.1 环 $=1\,\mathrm{mm}$，则弹着点定位误差应控制在约 $\pm0.5\,\mathrm{mm}$ 内。按“每个目标特征由 2～3 个像素采样”的经验，单像素对应的物方尺寸应 $\leq0.3\sim0.5\,\mathrm{mm}$。若需覆盖的靶面区域为 $L\times L$（mm），所需像素数 $N\geq(L/0.3)^2$。例如覆盖 $100\,\mathrm{mm}\times100\,\mathrm{mm}$ 时，至少约 $334\times334\approx11$ 万像素，实际选百万像素相机就很充裕。

② \textbf{实时判断中靶、不每枪换靶纸：}对靶面区域连续采集，用\textbf{帧间差分（图像相减）$+$ 阈值}检测新出现的暗斑（弹孔），触发识别程序计算环值；靶纸上的历史弹孔通过“当前帧与历史背景模型比较”来区分，从而支持多枪累积而不换纸。

③ \textbf{机位导致的图像变形：}相机与靶面不垂直会造成\textbf{透视变形}，使圆形靶环变成椭圆。需先做\textbf{相机标定}，用已知靶环（同心圆）做基准，估计单应/透视变换并做\textbf{几何校正}，把图像反变换到正视图后再计算。

④ \textbf{弹着点到靶心的最近距离：}弹孔是一小块区域，应计算其\textbf{轮廓上各点到靶心距离的最小值}（距离变换/逐点求最小），而不是弹孔中心到靶心的距离——这与射击规则中“以弹孔内沿到环线最近距离判定”一致。

\subsection*{G2\quad （隧道病害检测）要能分辨 $0.1\,\mathrm{mm}$ 的裂缝，检测车时速 $30\,\mathrm{km/h}$。请从分辨率与采样速度两方面说明对摄像机的需求。}
\textbf{横向分辨率：}若要分辨 0.1 mm，单像素物方尺寸应 $\leq0.05\,\mathrm{mm}$（留 2 倍裕量）。若横向视场为 $W$（如环向 3 m），则横向像素数 $N\geq W/0.05\,\mathrm{mm}=3000\,\mathrm{mm}/0.05\,\mathrm{mm}=60000$ ——单相机难以达到，需\textbf{多相机拼接}或提高物距精度。

\textbf{采样速度：}车速 $30\,\mathrm{km/h}=8.33\,\mathrm{m/s}$。若纵向也要 0.1 mm 以上的分辨能力（如 0.1 mm/行），行频需 $8330\,\mathrm{mm/s}\div0.1\,\mathrm{mm}\approx83300$ 行/s $\approx83\,\mathrm{kHz}$，远超常规相机，因此实践中往往：①用\textbf{线阵相机}配合运动扫描；②纵向适当放宽（裂缝是长条形，纵向精度可低），只保证横向 0.1 mm；③采用\textbf{多相机+高速传输/存储}。

\textbf{附加问题：}圆弧形表会导致\textbf{光照不均与几何畸变}，需\textbf{匀光增强}（非线性流形映射 / 自适应匀光）与几何校正；相机间图像有偏差时先做\textbf{融合与增强}再拼接。

\subsection*{G3\quad （方案选择）对一张低对比度车牌图像，要在 HE、CLAHE、伽马变换中选一种用于后续字符识别。你选哪个？说明理由。}
\textbf{答案：}优先选 \textbf{CLAHE}。理由：①车牌区域小、且常受局部光照不均影响，全局 HE 会把整幅图的暗部噪声一起放大、可能使车牌两侧过曝，反而降低识别率；②CLAHE 采用\textbf{分块 $+$ 限制对比度（clip limit）$+$ 双线性插值}，能增强局部对比度同时抑制噪声放大与块边界，非常适合后续字符识别；③伽马变换是全局点运算，只能整体提亮/压暗，对局部对比度不足无能为力，可作为预处理补充。\textbf{结论：主方案用 CLAHE，必要时串联 $\gamma<1$ 提亮暗部。}

\section*{题型 H\quad 选择／判断小题（快问快答）}
\begin{enumerate}[leftmargin=*,label=\arabic*.]
\item 下列哪个模板可用于去除椒盐噪声？\quad A. $3\times3$ 均值\quad B. 中值滤波\quad C. 拉普拉斯\quad D. 高提升\quad $\to$ \textbf{B}
\item 下列属于空间域点处理的是？\quad A. 均值滤波\quad B. 中值滤波\quad C. 伽马变换\quad D. 拉普拉斯锐化\quad $\to$ \textbf{C}（点运算只与单像素有关）
\item 对灰度阶跃最能产生“双响应”的是？\quad A. 一阶微分\quad B. 二阶微分\quad C. 均值滤波\quad D. 中值滤波\quad $\to$ \textbf{B}
\item 能同时“平滑噪声且较少模糊边缘”的是？\quad A. 均值\quad B. 高斯低通\quad C. 中值\quad D. 理想低通\quad $\to$ \textbf{C}
\item 8 bit 灰度图像最多能表示多少级灰度？\quad $\to$ \textbf{$2^8=256$ 级}
\item 直方图均衡的变换函数是什么？\quad $\to$ \textbf{累积分布函数（CDF）}
\item “同时对比”说明感知亮度依赖于？\quad $\to$ \textbf{周围（背景）亮度}
\item $m$ 邻接主要用于解决什么问题？\quad $\to$ \textbf{8 邻接路径/连通的二义性}
\item 纯锐化（高通）掩模的系数和是多少？\quad $\to$ \textbf{0}
\item CLAHE 相比 AHE 的关键改进是？\quad $\to$ \textbf{直方图裁剪（clip limit）$+$ 双线性插值}
\item 频域滤波前把图像乘 $(-1)^{x+y}$ 的作用？\quad $\to$ \textbf{把频谱直流分量移到中心（等价 fftshift）}
\item 下列最易产生振铃的是？\quad A. 高斯低通\quad B. 巴特沃斯低通($n=2$)\quad C. 理想低通\quad D. 均值滤波\quad $\to$ \textbf{C}
\item 去除图像上周期性条纹干扰宜用？\quad $\to$ \textbf{频域带阻／陷波滤波}
\end{enumerate}

\medskip
\textbf{最后一句：}小测题型即使变了，考查的“内核”仍是——\textbf{看系数和判滤波性质、看直方图判明暗对比、看微分阶数判边缘响应、看截止频率判通阻、看振铃判滤波器类型、按定义说清概念。}把第 6 部分 A/B/D/E 四类题刷熟，基本就稳了。

\begin{center}\small
高级图像处理 · 小测复习讲义 ｜ 依据 Lecture01 / 02-1 / 02-2、传统图像增强、高级图像处理02、Enhancement3 及课堂练习1 整理
\end{center}

\end{document}
